\section{Introduction}
\label{sec:introduction}

Autonomous multi-agent systems are designed to act at machine speed and scale without requiring human approval at every decision point. This capability is their core value proposition — and precisely the property that makes them structurally ungovernable without an architectural fallback. When an agent spawns a sub-agent, which spawns another, which spawns a fourth, the chain of causal responsibility for any resulting outcome becomes distributed across machine actors, none of whom individually hold terminal authority, and none of whom can be compelled to surrender it. We call this fallback the \emph{Bailout Protocol}: a compulsory exception-escalation mechanism embedded in the \bxthree{} Framework at the architectural level, not as a best practice or configurable safety option.

\subsection{Why Autonomous Systems Require a Mandatory Bailout Mechanism}

Multi-agent autonomous systems create two compounding accountability problems that no amount of logging, auditing, or policy guidance can solve architecturally. The first is the \emph{propagation-of-effects problem}. When an agent delegates a subtask to a spawned node, it delegates a portion of its own authority. The spawned node may in turn delegate further, creating a chain of sub-delegations that expands the system's effective operational scope beyond what any principal approved. Each link in this chain is, individually, rational — the delegating node was operating within its provisioned bounds, the sub-node was operating within the scope it was given. But at no point along the chain is there a mandatory mechanism that forces the exception upward to a human who can adjudicate whether the cumulative effect of all those individually-bounded actions is desirable. The propagation-of-effects problem is structural, not behavioral. It cannot be solved by instructing agents to escalate more frequently.

The second is the \emph{authority-gap problem}. In conventional software systems, authority is managed through access controls, review gates, and deployment pipelines — mechanisms enforced at design time and verified at deployment time. In autonomous multi-agent systems, agents act at runtime, often in response to inputs not anticipated at design time, and their actions may have consequences that propagate before any human has had the opportunity to evaluate them. The authority gap is the difference between the authority an agent was given at provisioning and the authority it effectively exercises during execution. As systems scale in depth and breadth, this gap widens. The result is a divergence between the system's actual authority and its authorized authority — a divergence that is invisible to the human operators who hold formal accountability for the outcomes.

Existing approaches to this problem fall into two categories, neither adequate. The first is \emph{post-hoc auditing}: after the fact, reconstruct what happened from logs and assign responsibility retrospectively. Post-hoc auditing is necessary but insufficient — it cannot prevent harm that has already propagated through the agent network. The second is \emph{runtime human-in-the-loop (HITL)}: require a human to approve each significant action before execution. HITL solves the authority gap but destroys operational utility — it requires the human to be present and attentive at a frequency that matches the agent's action rate, which is precisely the property that makes autonomous agents valuable in the first place. Both approaches treat the bailout path as an afterthought: something an agent \emph{can} do if it decides to, rather than something it \emph{must} do as a structural compulsion.

The \bxthree{} Framework takes a different approach. Rather than asking how humans can monitor agents more effectively, it asks what architectural conditions must hold so that accountability is preserved regardless of what agents do. The answer is that the accountability chain must be structurally non-negotiable: every spawned node must carry an immutable reference to a named human principal, embedded in the node's execution context in a way that cannot be overwritten by any machine actor. The \purpose{} layer carries intent and judgment. The \bounds{} layer enforces that no node may operate outside its provisioned operating range without triggering mandatory escalation. The \fact{} layer records every protocol event in an append-only forensic ledger, making accountability verifiable after the fact and enforceable before it. The Bailout Protocol is the mechanism that operationalizes this return to the human — compulsorily, deterministically, and with full forensic attribution at every step.

\subsection{What Happens When Systems Do Not Have One}

Without a mandatory bailout mechanism, multi-agent systems that spawn other agents exhibit a failure mode we term \emph{autonomous accountability collapse}. In this mode, exception states propagate through the machine-only execution graph indefinitely — absorbed by intermediate nodes, re-spawned as new tasks, or suppressed in the interest of maintaining operational continuity — and no human principal is ever notified, never mind given the opportunity to intervene. The system does not fail by stopping. It fails by continuing to act on the basis of a decision no human ever made.

This failure mode is not hypothetical. It is the expected behavior of any multi-agent system whose escalation logic is implemented as a configurable policy rather than a structural requirement. When escalation is a policy, it can be overridden by a higher-priority policy. When it is a configurable threshold, it can be raised or lowered. When it is a best practice, it can be deemed inapplicable in time-critical situations. Each relaxation of the requirement is individually rational. The cumulative effect is that the bailout path, when it is needed most, is not available.

The deeper problem is that autonomous accountability collapse is not detectable from within the system. A system that has lost its accountability anchor looks, from the inside, entirely functional: agents are spawning, tasks are being decomposed, actions are being executed. The collapse is visible only from the outside — in outcomes that diverge from authorized behavior, in audit trails that terminate at machine actors, in principals who discover after the fact that they bear accountability for decisions they did not know were being made. By the time the divergence becomes visible externally, the harm has already propagated.

\subsection{The Core Insight: Accountability Requires an Architectural Fallback}

The Bailout Protocol rests on a single architectural insight: \emph{accountability is not a property of agent behavior — it is a property of system architecture}. Accountability cannot be achieved by prescribing correct behavior to agents, because agents that spawn other agents inevitably expand their effective authority beyond what any prescription can cover. Accountability cannot be achieved by monitoring agent outputs, because outputs are a lagging indicator of a process that may already have caused harm. Accountability can only be achieved by making the chain of authority structurally non-negotiable — such that the human principal is not one possible responder to an exception but the only possible terminus for any exception that cannot be resolved within the machine-only execution graph.

This insight has three structural consequences for the design of the \bxthree{} Framework. First, the \purpose{} layer must encode the human principal's intent in a form that is immutable at the \fact{} layer and passed forward at spawn time. Every spawned node carries an immutable parent pointer $\alpha(n)$ that chains upward through the hierarchy and terminates at a human accountability anchor — established at provisioning and never modified for the node's operational lifetime. Second, the bailout path must be compulsive rather than elective. When a node's \bounds{} engine encounters an exception condition that exceeds its provisioned operating range, the node does not choose whether to escalate. It is structurally compelled to escalate. The decision is not made by the node — it is made by the architecture. Third, the forensic ledger maintained by the \fact{} layer must record every protocol event in an append-only chain that is tamper-evident and human-anchored. No machine actor can modify the ledger to obscure the chain of custody. No machine actor can issue a directive that appears, in the ledger, to have been issued by a human.

These three structural properties — immutable human anchoring, compulsive escalation, and forensic non-repudiation — are jointly necessary and jointly sufficient for the \emph{Bailout Invariant}: for any exception event in any \bxthree{} node, there exists exactly one identifiable human principal who bears accountability for the outcome, and that principal is reachable within a bounded, finite number of propagation steps. The invariant holds unconditionally of machine-actor behavior — regardless of how agents behave, what exceptions arise, or how deeply the hierarchy is nested.

\subsection{Paper Roadmap}

The remainder of this paper is organized as follows. Section~\ref{sec:background} reviews the relevant literature on accountability in multi-agent AI systems, human-in-the-loop design, and fail-safe design in safety-critical systems, and situates the \bxthree{} Framework's three-layer model within this tradition. Section~\ref{sec:bailout} specifies the Bailout Protocol as a deterministic state machine, defines the trigger predicate that activates it, describes the escalation algorithm that propagates exceptions upward through the immutable parent chain, and details the bypass mechanism that prevents machine actors from serving as terminal resolvers. Section~\ref{sec:bailout-trigger} characterizes the exception conditions that trigger bailout and the sealed-envelope architecture that prevents trigger suppression. Section~\ref{sec:correctness} states the three correctness properties — Safety, Liveness, and Accountability — provides proof sketches for each, and derives the Bailout Invariant as their logical consequence. Section~\ref{sec:limitations} identifies four honest limitations: human response latency, adversarial hardening of trigger predicates, scalability under one-to-one anchoring, and exploitation risk from misconfigured bounds. Section~\ref{sec:future-work} outlines three directions for future research: formal verification of core theorems, adaptive timeout and priority escalation, and distributed human anchor pools for industrial-scale deployment. Section~\ref{sec:conclusion} summarizes the contributions and argues that the Bailout Protocol, as a structural instantiation of the \bxthree{} Framework's three-layer architecture, is not merely compatible with autonomous multi-agent operation but is its necessary architectural precondition.
